\makeabstract
    {
    Microwaving Molecules Until They Tell Us Their Secrets: Spectroscopic Characterization of a Halopropene
    }
    {
    Noah Keil\textsuperscript{1}, Asher Navas\textsuperscript{1}, Mark D. Marshall\textsuperscript{1}, Helen O. Leung\textsuperscript{1}
    }
    {
    \textsuperscript{1} Amherst College Dept. of Chemistry
    }
    {
    In quantum systems such as molecules and their complexes, angular momentum is quantized. This quantization yields distinct transitions that can be measured and exploited to determine the precise structure of the molecule and its complexes with other atoms. By determining the structure of such a system, a fuller understanding of intermolecular forces can be advanced. In this project, we use microwave spectroscopy to characterize the compound 3-chloro-1,1,3,3-tetrafluropropene and its complex with argon.
    }
    {
    An experimental spectrum was measured using Chirped-Pulse Fourier-Transform Rotational Microwave (CP-FTMW) Spectroscopy. To predict discrete transitions, the system was modeled with the quantum chemistry software Gaussian16. Subsequently, the theoretical spectrum was simulated and assigned to the experimental spectrum using PGOPHER, and revised structural coordinates were calculated where possible. 
    }
    {
    A dihedral scan of the monomer revealed two enantiomeric rotamers, which have identical spectroscopic constants. We then matched peaks for the Cl35 and Cl37 isotopologues of the monomer. We fit rotational constants, distortion constants, and quadrupole coupling constants to the best of their ability based on the transitions assigned. We then scanned Argon around the monomer for the argon complex, and we took the lowest energy configuration to assign transitions for the Cl35 and Cl37 isotopologues of the argon complex.
    }
    {
    The theory we used provided an estimate that lay within 10\% of the value of our assigned constants, producing a relatively accurate theoretical structure. In the future, we want to look at the acetylene complex with the monomer to find out more about the structure of the molecule and its intermolecular interactions.
    }

    \newpage
    